% ch01_introduction.tex — Introduction and Motivations
\chapter{Introduction and Motivations}\label{ch:introduction}

Consider a barista stirring a cocktail in a glass. After vigorously swirling
the liquid, she can be certain that at least one point in the liquid ends up
exactly at its original position. This result, known as Brouwer's theorem,
illustrates a simple but profound idea: under certain conditions, every
transformation has at least one \emph{fixed point}, a point that does not
move. This idea runs through mathematics: it appears in analysis, topology,
algebra, economics (Nash equilibrium), computer science (program semantics),
and game theory. This course systematically explores fixed-point theorems,
their proofs, and their applications.

% ------------------------------------------------------------
\section{What is a fixed point?}
% ------------------------------------------------------------

The definition is disarmingly simple --- but do not be fooled.

\begin{definition}[Fixed point]
Let $X$ be a set and $T : X \to X$ a mapping. An element $x^* \in X$ is a
\emph{fixed point} of $T$ if $T(x^*) = x^*$. The set of fixed points of $T$
is denoted
\[
    \mathrm{Fix}(T) = \{ x \in X : T(x) = x \}.
\]
\end{definition}

This seemingly elementary definition leads to remarkably deep questions as
soon as $X$ is endowed with additional structure (metric, topology, order) and
conditions are imposed on $T$.

\begin{remark}
The equation $T(x) = x$ can be reformulated as $F(x) = 0$ by setting
$F = T - \mathrm{Id}$. Conversely, any equation $F(x) = 0$ can be reduced to a
fixed point problem by setting $T(x) = x - F(x)$ (or $T(x) = x + \alpha F(x)$
for a suitably chosen parameter $\alpha$). This duality underlies many
numerical methods.
\end{remark}

\section{Historical context}

Fixed point theory has a rich and fascinating history.

\subsection{Origins: Picard and Banach}

The method of successive approximations dates back to the work of \'Emile Picard
(1890) on ordinary differential equations. Picard showed that the equation
\[
    y'(t) = f(t, y(t)), \quad y(t_0) = y_0
\]
has a unique solution under a Lipschitz condition on $f$ with respect to $y$,
by constructing the sequence of iterates
\[
    y_{n+1}(t) = y_0 + \int_{t_0}^{t} f(s, y_n(s)) \, ds.
\]

It was Stefan Banach who, in his doctoral thesis (1920, published 1922),
formalized this idea by stating the \emph{Contraction Principle} in the general
framework of complete metric spaces.

\subsection{Brouwer's theorem}

In 1911, L.\,E.\,J.\ Brouwer proved that in finite dimensions, every continuous
map from a compact convex set into itself has a fixed point. This result, purely
topological in nature (no metric hypothesis), relies on arguments from
topological degree theory and algebraic topology.

\subsection{Infinite-dimensional extensions}

Juliusz Schauder (1930) generalized Brouwer's theorem to infinite-dimensional
Banach spaces by replacing compactness of the set with compactness of the
operator. Shizuo Kakutani (1941) extended the result to multivalued
correspondences (set-valued mappings).

\subsection{Order-theoretic approach}

Alfred Tarski (1955) established a fixed point theorem in complete lattices,
based solely on monotonicity of the mapping, with no metric or topological
hypothesis. This result, often called the Knaster--Tarski theorem, finds
applications in theoretical computer science and logic.

\section{Elementary examples}

\begin{example}[Fixed point in dimension 1]
Let $f : [0, 1] \to [0, 1]$ be a continuous map. By the Intermediate Value
Theorem, $f$ has a fixed point.

Indeed, set $g(x) = f(x) - x$. Then $g(0) = f(0) \geq 0$ and
$g(1) = f(1) - 1 \leq 0$. By the Intermediate Value Theorem, there exists
$x^* \in [0,1]$ such that $g(x^*) = 0$, i.e., $f(x^*) = x^*$.
\end{example}

\begin{example}[Non-existence of fixed points]
The map $T : \R \to \R$ defined by $T(x) = x + 1$ has no fixed point. This
shows the necessity of working in a ``sufficiently small'' set (bounded,
compact, etc.).
\end{example}

\begin{example}[Multiplicity of fixed points]
The identity map $T(x) = x$ has every point as a fixed point. The map
$T(x) = x^2$ on $\R$ has exactly two fixed points: $0$ and $1$.
\end{example}

\begin{example}[Rotation without fixed point]
The rotation $R_\theta : S^1 \to S^1$ by angle $\theta \neq 0 \pmod{2\pi}$ on
the unit circle has no fixed point. This shows that convexity is essential in
Brouwer's theorem: the circle $S^1$ is compact but not convex.
\end{example}

\section{Formulation as a functional equation}

Many classical mathematical problems can be reformulated as fixed point problems.

\subsection{Integral equations}

The Fredholm integral equation of the second kind
\[
    u(t) = f(t) + \lambda \int_a^b K(t,s) \, u(s) \, ds
\]
is of the form $u = Tu$ where the operator $T$ is defined by
\[
    (Tu)(t) = f(t) + \lambda \int_a^b K(t,s) \, u(s) \, ds.
\]

\subsection{Differential equations}

The Cauchy problem $y' = f(t,y)$, $y(t_0) = y_0$ is equivalent, under
regularity assumptions, to the fixed point problem $y = Ty$ where
\[
    (Ty)(t) = y_0 + \int_{t_0}^t f(s, y(s)) \, ds.
\]

\subsection{Partial differential equations}

Boundary value problems for elliptic PDEs are often reformulated as fixed point
problems in Sobolev spaces. For example, the problem
\[
    -\Delta u = g(u) \text{ in } \Omega, \quad u = 0 \text{ on } \partial\Omega
\]
becomes $u = (-\Delta)^{-1} g(u) = T(u)$.

\subsection{Nash equilibria}

A Nash equilibrium in an $n$-player game can be viewed as a fixed point of a
best-response correspondence. If $S_i$ is the strategy set of player $i$ and
$BR_i(s_{-i})$ is their best response to the other players' strategies, then a
Nash equilibrium is a profile $s^* = (s_1^*, \ldots, s_n^*)$ such that
$s_i^* \in BR_i(s_{-i}^*)$ for all $i$.

\section{Three families of theorems}

Fixed point theorems fall into three broad families, according to the
mathematical structure employed.

\begin{intuition}[title=\textbf{The three approaches}]
\begin{enumerate}
    \item \textbf{Metric approach}: one assumes that $T$ contracts distances.
          The fixed point is unique and obtained by iteration. \\
          Prototype: Banach's theorem.
    \item \textbf{Topological approach}: one assumes continuity of $T$ and
          compactness of the domain. Existence is guaranteed, but neither
          uniqueness nor constructibility. \\
          Prototype: Brouwer's theorem.
    \item \textbf{Order-theoretic approach}: one assumes monotonicity of $T$ in
          a complete lattice. Existence is guaranteed, and the set of fixed
          points itself forms a complete lattice. \\
          Prototype: Tarski's theorem.
\end{enumerate}
\end{intuition}

\section{Preliminaries and background}

\subsection{Metric spaces}

\begin{definition}[Metric space]
A \emph{metric space} is a pair $(X, d)$ where $X$ is a set and
$d : X \times X \to [0, +\infty)$ satisfies for all $x, y, z \in X$:
\begin{enumerate}
    \item $d(x, y) = 0 \Leftrightarrow x = y$ \quad (separation),
    \item $d(x, y) = d(y, x)$ \quad (symmetry),
    \item $d(x, z) \leq d(x, y) + d(y, z)$ \quad (triangle inequality).
\end{enumerate}
\end{definition}

\begin{definition}[Completeness]
A metric space $(X, d)$ is \emph{complete} if every Cauchy sequence in $X$
converges. Recall that a sequence $(x_n)$ is Cauchy if for every
$\varepsilon > 0$, there exists $N \in \N$ such that $d(x_m, x_n) < \varepsilon$
for all $m, n \geq N$.
\end{definition}

\begin{definition}[Lipschitz mapping]
Let $(X, d_X)$ and $(Y, d_Y)$ be metric spaces. A mapping $T : X \to Y$ is
\emph{Lipschitz} with constant $L \geq 0$ if
\[
    d_Y(T(x), T(y)) \leq L \, d_X(x, y) \quad \forall x, y \in X.
\]
If $L < 1$, $T$ is called a \emph{contraction}. If $L = 1$, $T$ is called
\emph{nonexpansive}.
\end{definition}

\subsection{Banach spaces}

\begin{definition}[Banach space]
A \emph{Banach space} is a normed vector space $(E, \norm{\cdot})$ that is
complete with respect to the induced metric $d(x,y) = \norm{x - y}$.
\end{definition}

\begin{definition}[Compactness]
A subset $K$ of a topological space is \emph{compact} if every open cover of
$K$ has a finite subcover. In a metric space, this is equivalent to sequential
compactness: every sequence in $K$ has a convergent subsequence with limit
in $K$.
\end{definition}

\begin{remark}
In infinite dimensions, closed bounded balls are never compact (Riesz's theorem).
This is why topological fixed point theorems in infinite dimensions (Schauder)
require compactness hypotheses on the operator rather than on the domain.
\end{remark}

\subsection{Convexity}

\begin{definition}[Convex set]
A subset $C$ of a vector space $E$ is \emph{convex} if for all $x, y \in C$
and all $t \in [0,1]$, we have $(1-t)x + ty \in C$.
\end{definition}

\begin{definition}[Convex hull]
The \emph{convex hull} of a set $A \subset E$ is the smallest convex set
containing $A$:
\[
    \mathrm{co}(A) = \left\{ \sum_{i=1}^n \lambda_i x_i : n \in \N^*,
    \, x_i \in A, \, \lambda_i \geq 0, \, \sum_{i=1}^n \lambda_i = 1 \right\}.
\]
\end{definition}

\section{Picard iteration}

The most elementary method for finding a fixed point is Picard iteration
(the method of successive approximations).

\begin{algorithme}[title=\textbf{Picard iteration}]
\textbf{Input}: Mapping $T : X \to X$, initial point $x_0 \in X$. \\
\textbf{Procedure}: Define the sequence $(x_n)$ by $x_{n+1} = T(x_n)$ for $n \geq 0$. \\
\textbf{Output}: If $(x_n)$ converges to $x^*$ and $T$ is continuous, then $x^*$ is a fixed point of $T$.
\end{algorithme}

\begin{remark}
Convergence of Picard iteration is not guaranteed in general. The Banach theorem
(Chapter~2) provides sufficient conditions.
\end{remark}

\section{The one-dimensional fixed point theorem}

\begin{theorem}[One-dimensional fixed point theorem]
Let $f : [a,b] \to [a,b]$ be a continuous map. Then $f$ has at least one fixed
point in $[a,b]$.
\end{theorem}

\begin{proof}
Consider the function $g : [a,b] \to \R$ defined by $g(x) = f(x) - x$. Since
$f([a,b]) \subset [a,b]$, we have $g(a) = f(a) - a \geq 0$ and
$g(b) = f(b) - b \leq 0$. By the Intermediate Value Theorem, there exists
$x^* \in [a,b]$ such that $g(x^*) = 0$, i.e., $f(x^*) = x^*$.
\end{proof}

\begin{warning}[title=\textbf{Necessary hypotheses}]
Both hypotheses are essential:
\begin{itemize}
    \item \textbf{Continuity}: the map $f : [0,1] \to [0,1]$ defined by
          $f(x) = 0$ if $x > 1/2$ and $f(x) = 1$ if $x \leq 1/2$ has no fixed point.
    \item \textbf{Invariance} ($f([a,b]) \subset [a,b]$): the map $f(x) = x + 1$
          on $[0,1]$ has no fixed point since $f([0,1]) = [1,2] \not\subset [0,1]$.
\end{itemize}
\end{warning}

\section{Topological degree: a preview}

The topological degree is an algebraic invariant that ``counts'' the solutions
of an equation $F(x) = p$ with multiplicity and sign. For a $C^1$ map
$F : \overline{\Omega} \to \R^n$, with $\Omega$ a bounded open subset of $\R^n$
and $p \notin F(\partial\Omega)$, the Brouwer degree is defined by
\[
    \deg(F, \Omega, p) = \sum_{x \in F^{-1}(p)} \mathrm{sgn}\,\det(DF(x)).
\]

\begin{theorem}[Existence via nonzero degree]
If $\deg(F, \Omega, p) \neq 0$, then the equation $F(x) = p$ has at least one
solution in $\Omega$.
\end{theorem}

This connection between algebraic topology and existence of solutions will be
developed in detail in Chapter~5.

\section{Course outline}

\begin{keyformulas}[title=\textbf{Course architecture}]
\begin{center}
\renewcommand{\arraystretch}{1.3}
\begin{tabular}{|c|l|l|}
\hline
\textbf{Ch.} & \textbf{Title} & \textbf{Main result} \\
\hline
1 & Introduction & Motivations, preliminaries \\
2 & Banach & Contraction Principle \\
3 & Metric extensions & Edelstein, Kannan, \'Ciri\'c, Caristi \\
4 & Generalized spaces & $b$-metrics, $G$-metrics \\
5 & Brouwer & Topological fixed point (finite dim.) \\
6 & Schauder & Topological fixed point (infinite dim.) \\
7 & Kakutani & Multivalued correspondences \\
8 & Ordered spaces & Tarski--Knaster \\
9 & Applications & ODE, PDE, games, economics \\
10 & Markov--Kakutani & Common fixed points \\
11 & Random fixed points & Measurable theory \\
\hline
\end{tabular}
\end{center}
\end{keyformulas}

\section{Exercises}

\begin{exercise}
Let $f : [0,1] \to [0,1]$ be an increasing (not necessarily continuous) map.
Show that $f$ has a fixed point.
\textit{Hint: consider $x^* = \sup\{x \in [0,1] : f(x) \geq x\}$.}
\end{exercise}

\begin{exercise}
Let $f : \R \to \R$ be continuous with $\abs{f(x) - f(y)} < \abs{x - y}$ for
all $x \neq y$. Show that $f$ has at most one fixed point. Give an example
where $f$ has no fixed point.
\end{exercise}

\begin{exercise}
Let $T : \R^n \to \R^n$ be an affine isometry (i.e.,
$\norm{T(x) - T(y)} = \norm{x - y}$ for all $x, y$). Show that either $T$ has
a fixed point or $T$ is a translation with no fixed point.
\end{exercise}

\begin{exercise}
Show that the rotation of $\R^2$ by angle $\pi/4$ about the origin has a unique
fixed point. More generally, what can be said about the fixed points of a
rotation by angle $\theta$ in dimension $n$?
\end{exercise}

\begin{exercise}
Let $(X, d)$ be a compact metric space and $T : X \to X$ an isometry (i.e.,
$d(T(x), T(y)) = d(x,y)$ for all $x,y$). Show that $T$ is surjective.
Deduce that if $X$ is moreover convex in a normed space, then $T$ has a fixed
point.
\end{exercise}

\begin{exercise}
Let $A$ be an $n \times n$ matrix with nonnegative entries such that the sum of
the entries in each column equals $1$ (stochastic matrix). Show that $A$ has an
eigenvector associated to the eigenvalue $1$ in the simplex
$\Delta_n = \{x \in \R^n : x_i \geq 0, \sum x_i = 1\}$. Interpret this result
as a fixed point theorem.
\end{exercise}
